% Options for packages loaded elsewhere
% Options for packages loaded elsewhere
\PassOptionsToPackage{unicode}{hyperref}
\PassOptionsToPackage{hyphens}{url}
\PassOptionsToPackage{dvipsnames,svgnames,x11names}{xcolor}
%
\documentclass[
  letterpaper,
  DIV=11,
  numbers=noendperiod]{scrartcl}
\usepackage{xcolor}
\usepackage{amsmath,amssymb}
\setcounter{secnumdepth}{-\maxdimen} % remove section numbering
\usepackage{iftex}
\ifPDFTeX
  \usepackage[T1]{fontenc}
  \usepackage[utf8]{inputenc}
  \usepackage{textcomp} % provide euro and other symbols
\else % if luatex or xetex
  \usepackage{unicode-math} % this also loads fontspec
  \defaultfontfeatures{Scale=MatchLowercase}
  \defaultfontfeatures[\rmfamily]{Ligatures=TeX,Scale=1}
\fi
\usepackage{lmodern}
\ifPDFTeX\else
  % xetex/luatex font selection
\fi
% Use upquote if available, for straight quotes in verbatim environments
\IfFileExists{upquote.sty}{\usepackage{upquote}}{}
\IfFileExists{microtype.sty}{% use microtype if available
  \usepackage[]{microtype}
  \UseMicrotypeSet[protrusion]{basicmath} % disable protrusion for tt fonts
}{}
\makeatletter
\@ifundefined{KOMAClassName}{% if non-KOMA class
  \IfFileExists{parskip.sty}{%
    \usepackage{parskip}
  }{% else
    \setlength{\parindent}{0pt}
    \setlength{\parskip}{6pt plus 2pt minus 1pt}}
}{% if KOMA class
  \KOMAoptions{parskip=half}}
\makeatother
% Make \paragraph and \subparagraph free-standing
\makeatletter
\ifx\paragraph\undefined\else
  \let\oldparagraph\paragraph
  \renewcommand{\paragraph}{
    \@ifstar
      \xxxParagraphStar
      \xxxParagraphNoStar
  }
  \newcommand{\xxxParagraphStar}[1]{\oldparagraph*{#1}\mbox{}}
  \newcommand{\xxxParagraphNoStar}[1]{\oldparagraph{#1}\mbox{}}
\fi
\ifx\subparagraph\undefined\else
  \let\oldsubparagraph\subparagraph
  \renewcommand{\subparagraph}{
    \@ifstar
      \xxxSubParagraphStar
      \xxxSubParagraphNoStar
  }
  \newcommand{\xxxSubParagraphStar}[1]{\oldsubparagraph*{#1}\mbox{}}
  \newcommand{\xxxSubParagraphNoStar}[1]{\oldsubparagraph{#1}\mbox{}}
\fi
\makeatother


\usepackage{longtable,booktabs,array}
\newcounter{none} % for unnumbered tables
\usepackage{calc} % for calculating minipage widths
% Correct order of tables after \paragraph or \subparagraph
\usepackage{etoolbox}
\makeatletter
\patchcmd\longtable{\par}{\if@noskipsec\mbox{}\fi\par}{}{}
\makeatother
% Allow footnotes in longtable head/foot
\IfFileExists{footnotehyper.sty}{\usepackage{footnotehyper}}{\usepackage{footnote}}
\makesavenoteenv{longtable}
\usepackage{graphicx}
\makeatletter
\newsavebox\pandoc@box
\newcommand*\pandocbounded[1]{% scales image to fit in text height/width
  \sbox\pandoc@box{#1}%
  \Gscale@div\@tempa{\textheight}{\dimexpr\ht\pandoc@box+\dp\pandoc@box\relax}%
  \Gscale@div\@tempb{\linewidth}{\wd\pandoc@box}%
  \ifdim\@tempb\p@<\@tempa\p@\let\@tempa\@tempb\fi% select the smaller of both
  \ifdim\@tempa\p@<\p@\scalebox{\@tempa}{\usebox\pandoc@box}%
  \else\usebox{\pandoc@box}%
  \fi%
}
% Set default figure placement to htbp
\def\fps@figure{htbp}
\makeatother





\setlength{\emergencystretch}{3em} % prevent overfull lines

\providecommand{\tightlist}{%
  \setlength{\itemsep}{0pt}\setlength{\parskip}{0pt}}



 


\KOMAoption{captions}{tableheading}
\makeatletter
\@ifpackageloaded{caption}{}{\usepackage{caption}}
\AtBeginDocument{%
\ifdefined\contentsname
  \renewcommand*\contentsname{Table of contents}
\else
  \newcommand\contentsname{Table of contents}
\fi
\ifdefined\listfigurename
  \renewcommand*\listfigurename{List of Figures}
\else
  \newcommand\listfigurename{List of Figures}
\fi
\ifdefined\listtablename
  \renewcommand*\listtablename{List of Tables}
\else
  \newcommand\listtablename{List of Tables}
\fi
\ifdefined\figurename
  \renewcommand*\figurename{Figure}
\else
  \newcommand\figurename{Figure}
\fi
\ifdefined\tablename
  \renewcommand*\tablename{Table}
\else
  \newcommand\tablename{Table}
\fi
}
\@ifpackageloaded{float}{}{\usepackage{float}}
\floatstyle{ruled}
\@ifundefined{c@chapter}{\newfloat{codelisting}{h}{lop}}{\newfloat{codelisting}{h}{lop}[chapter]}
\floatname{codelisting}{Listing}
\newcommand*\listoflistings{\listof{codelisting}{List of Listings}}
\makeatother
\makeatletter
\makeatother
\makeatletter
\@ifpackageloaded{caption}{}{\usepackage{caption}}
\@ifpackageloaded{subcaption}{}{\usepackage{subcaption}}
\makeatother
\usepackage{bookmark}
\IfFileExists{xurl.sty}{\usepackage{xurl}}{} % add URL line breaks if available
\urlstyle{same}
\hypersetup{
  pdftitle={Lembar Jawaban: Persiapan UTS Probabilitas dan Statistika 2021},
  colorlinks=true,
  linkcolor={blue},
  filecolor={Maroon},
  citecolor={Blue},
  urlcolor={Blue},
  pdfcreator={LaTeX via pandoc}}


\title{Lembar Jawaban: Persiapan UTS Probabilitas dan Statistika 2021}
\author{}
\date{}
\begin{document}
\maketitle


\textbf{Nama Dosen:} Dimitri Mahayana (Asumsi)\\
\textbf{Mata Kuliah:} Probabilitas dan Statistika

\begin{center}\rule{0.5\linewidth}{0.5pt}\end{center}

\subsubsection{\texorpdfstring{\textbf{Soal 1: Teorema Bayes (Diagnosa
Penyakit)}}{Soal 1: Teorema Bayes (Diagnosa Penyakit)}}\label{soal-1-teorema-bayes-diagnosa-penyakit}

\textbf{Soal:} Penyakit langka menyerang 1 dari 1.000 orang. Tes
mendeteksi hasil positif 99\% benar jika sakit, namun ada 2\% false
positive (salah deteksi positif pada orang sehat). Jika seseorang dites
dan hasilnya positif, berapa peluang ia sungguh-sungguh terinfeksi?
Berikan analisis dan saran!

\textbf{Jawaban:} Diketahui probabilitas awal (Prior): * Peluang sakit
\(P(D) = 1/1000 = 0,001\) * Peluang sehat \(P(D') = 1 - 0,001 = 0,999\)
Probabilitas bersyarat (Likelihood): * Positif jika sakit (\emph{True
Positive}) \(P(+|D) = 0,99\) * Positif jika sehat (\emph{False
Positive}) \(P(+|D') = 0,02\)

Berdasarkan \textbf{Teorema Bayes}, probabilitas seseorang benar-benar
sakit jika tesnya positif adalah:
\[ P(D|+) = \frac{P(+|D) \cdot P(D)}{P(+|D) \cdot P(D) + P(+|D') \cdot P(D')} \]
\[ P(D|+) = \frac{(0,99)(0,001)}{(0,99)(0,001) + (0,02)(0,999)} \]
\[ P(D|+) = \frac{0,00099}{0,00099 + 0,01998} = \frac{0,00099}{0,02097} \approx \mathbf{0,0472} \text{ atau } \mathbf{4,72\%} \]

\textbf{Analisis \& Saran:} * \textbf{Analisis:} Peluang orang tersebut
benar-benar sakit meskipun tesnya positif ternyata sangat rendah (hanya
\textasciitilde4,72\%). Hal ini terjadi karena penyakit tersebut sangat
langka (1:1000) sehingga jumlah \emph{false positive} dari populasi yang
sehat jauh lebih banyak daripada jumlah orang sakit yang sebenarnya
(dikenal sebagai \emph{Base Rate Fallacy}). * \textbf{Saran:} Mengingat
persoalan ini menyangkut hidup dan mati, saran terbaik adalah
\textbf{tidak perlu panik terlebih dahulu} dan \textbf{lakukan tes
kedua} (dengan metode atau alat uji yang lebih spesifik) untuk
mengonfirmasi hasil positif pertama, sebelum mengambil tindakan medis
ekstrem.

\begin{center}\rule{0.5\linewidth}{0.5pt}\end{center}

\subsubsection{\texorpdfstring{\textbf{Soal 2: Probabilitas Empiris dari
Tabel
Frekuensi}}{Soal 2: Probabilitas Empiris dari Tabel Frekuensi}}\label{soal-2-probabilitas-empiris-dari-tabel-frekuensi}

\textbf{Soal:} a. Kecepatan downstream (download) harus minimal 500 kbps
(0,5 Mbps). Tentukan probabilitas TV internet bisa dinikmati dengan
baik! b. Upload file video 63 MB waktu tidak lebih dari 21 menit (1 Bps
= 8 bps). Tentukan probabilitasnya!

\textbf{Jawaban:} \textbf{a. Probabilitas Downstream \(\ge\) 500 kbps
(0,5 Mbps):} Dari tabel ``Kecepatan Download'', kelas yang bernilai
\(< 0,5\) Mbps adalah kelas No.~1 hingga 5.
\[ P(< 0,5 \text{ Mbps}) = 6,53\% + 33,42\% + 3,77\% + 8,79\% + 1,76\% = 54,27\% \]
Maka, probabilitas kecepatannya memadai (\(\ge 0,5\) Mbps) adalah:
\[ P(\ge 0,5 \text{ Mbps}) = 100\% - 54,27\% = \mathbf{45,73\%} \]

\textbf{b. Probabilitas Upload 63 MB dalam \(\le\) 21 menit:} * Konversi
ukuran file ke bit:
\(63 \text{ MB} = 63 \times 10^6 \text{ Bytes} \times 8 \text{ bits/Byte} = 504 \times 10^6 \text{ bits} = 504 \text{ Mbits}\).
* Konversi waktu maksimal:
\(21 \text{ menit} = 21 \times 60 \text{ detik} = 1260 \text{ detik}\).
* Kebutuhan kecepatan minimum:
\(V = \frac{504 \text{ Mbits}}{1260 \text{ detik}} = \mathbf{0,4 \text{ Mbps}}\).

Dari tabel ``Kecepatan Upload'', kelas yang bernilai \(< 0,4\) Mbps
adalah kelas No.~1 hingga 4.
\[ P(< 0,4 \text{ Mbps}) = 14,82\% + 13,32\% + 12,31\% + 13,82\% = 54,27\% \]
Maka, probabilitas upload selesai tepat waktu (\(\ge 0,4\) Mbps) adalah:
\[ P(\ge 0,4 \text{ Mbps}) = 100\% - 54,27\% = \mathbf{45,73\%} \]

\begin{center}\rule{0.5\linewidth}{0.5pt}\end{center}

\subsubsection{\texorpdfstring{\textbf{Soal 3: Distribusi
Pareto}}{Soal 3: Distribusi Pareto}}\label{soal-3-distribusi-pareto}

\textbf{Soal:} Running time \(R\) memiliki peluang \(R > 10\) adalah
\(0,5\). Tentukan peluang \(R > 1000\) unit waktu!
\(F(r) = (1 - r^{-a}) u(r-1)\)

\textbf{Jawaban:} Diketahui probabilitas pelengkapnya:
\(P(R > r) = 1 - F(r) = r^{-a}\). Kita cari nilai parameter \(a\)
menggunakan data awal: \[ P(R > 10) = 10^{-a} = 0,5 \] Dengan persamaan
eksponensial ini, kita bisa mensubstitusikannya langsung untuk mencari
target \(P(R > 1000)\): \[ P(R > 1000) = 1000^{-a} \] Karena
\(1000 = 10^3\), maka persamaannya dapat diubah menjadi:
\[ P(R > 1000) = (10^3)^{-a} = (10^{-a})^3 \]
\[ P(R > 1000) = (0,5)^3 = \mathbf{0,125} \]

\begin{center}\rule{0.5\linewidth}{0.5pt}\end{center}

\subsubsection{\texorpdfstring{\textbf{Soal 4: Distribusi Binomial
(Sistem Data
Center)}}{Soal 4: Distribusi Binomial (Sistem Data Center)}}\label{soal-4-distribusi-binomial-sistem-data-center}

\textbf{Soal:} 3 set infrastruktur di 3 Data Center independen. Peluang
bekerja dengan baik \(p = 0,9\). Definisikan \(X\) = jumlah Data Center
yang hidup. a. Tabel distribusi probabilitas b. Gambarkan histogram
\emph{(Secara Deskriptif)} c.~Mean dan Variansi d.~Sistem down total
(\(X=0\)) e. Sistem nyaris down, sisa 1 (\(X=1\)) f.~Minimal 1 DC
bekerja (\(X \ge 1\))

\textbf{Jawaban:} Variabel acak
\(X \sim \text{Binomial}(n=3, p=0,9, q=0,1)\). Rumus:
\(P(X=x) = \binom{3}{x} (0,9)^x (0,1)^{3-x}\)

\textbf{a. Tabel Distribusi Probabilitas:} \textbar{} \(x\) \textbar{}
\(P(X=x)\) \textbar{} Perhitungan \textbar{}
\textbar:---:\textbar:---\textbar:---\textbar{} \textbar{} 0 \textbar{}
0,001 \textbar{} \(\binom{3}{0} (0,9)^0 (0,1)^3\) \textbar{} \textbar{}
1 \textbar{} 0,027 \textbar{} \(\binom{3}{1} (0,9)^1 (0,1)^2\)
\textbar{} \textbar{} 2 \textbar{} 0,243 \textbar{}
\(\binom{3}{2} (0,9)^2 (0,1)^1\) \textbar{} \textbar{} 3 \textbar{}
0,729 \textbar{} \(\binom{3}{3} (0,9)^3 (0,1)^0\) \textbar{}

\textbf{b. Histogram:} Sumbu X melambangkan jumlah DC (0, 1, 2, 3) dan
Sumbu Y melambangkan probabilitas. Grafiknya akan sangat condong ke
kanan (miring negatif), dengan puncak batang tertinggi berada di \(X=3\)
sebesar 0,729 dan batang terpendek di \(X=0\) sebesar 0,001.

\textbf{c.~Mean dan Variansi:} * Mean
\(\mu = n \cdot p = 3(0,9) = \mathbf{2,7}\) * Variansi
\(\sigma^2 = n \cdot p \cdot q = 3(0,9)(0,1) = \mathbf{0,27}\)

\textbf{d, e, f.~Probabilitas Skenario:} * \textbf{d.} Probabilitas down
total (\(X=0\)) = \(\mathbf{0,001}\) * \textbf{e.} Probabilitas nyaris
down / 1 DC hidup (\(X=1\)) = \(\mathbf{0,027}\) * \textbf{f.}
Probabilitas minimal 1 hidup (\(X \ge 1\)) =
\(1 - P(X=0) = 1 - 0,001 = \mathbf{0,999}\)

\begin{center}\rule{0.5\linewidth}{0.5pt}\end{center}

\subsubsection{\texorpdfstring{\textbf{Soal 5: Keandalan Jaringan
Komunikasi
(Reliability)}}{Soal 5: Keandalan Jaringan Komunikasi (Reliability)}}\label{soal-5-keandalan-jaringan-komunikasi-reliability}

\textbf{Soal:} Probabilitas link baik \(p = 0,6\). Terdapat 8 link
(a,b,c,d,e,f,g,h). a) i. Tepat 2 link bekerja? ii. Link g dan tepat satu
link lagi bekerja? b) Jika tepat 6 link mati (berarti 2 hidup), peluang
A bisa berkomunikasi dengan B?

\textbf{Jawaban:} Berdasarkan deskripsi gambar standar untuk model
bridge (jembatan) dengan by-pass: Asumsikan total \(n=8\) link
independen.

\textbf{a. Pengamatan Waktu Acak (Distribusi Binomial n=8, p=0,6):} *
\textbf{i.} Tepat 2 link bekerja (\(X=2\)):
\[ P(X=2) = \binom{8}{2} (0,6)^2 (0,4)^6 = 28 \times 0,36 \times 0,004096 \approx \mathbf{0,0413} \]
* \textbf{ii.} Link \textbf{g} dan tepat \textbf{satu link lainnya}
bekerja: Karena link \textbf{g} dipastikan bekerja (probabilitas = 0,6),
maka dari sisa 7 link lainnya, hanya 1 yang boleh bekerja.
\[ P(\text{g hidup}) \times P(Y=1 \text{ dari } 7) = 0,6 \times \left[ \binom{7}{1} (0,6)^1 (0,4)^6 \right] \]
\[ = 0,6 \times [ 7 \times 0,6 \times 0,004096 ] = \mathbf{0,01032} \]

\textbf{b. Probabilitas Komunikasi A ke B jika tepat 2 link hidup:}
Berdasarkan diagram \emph{bridge network}, agar terminal A dapat
terhubung ke terminal B menggunakan \textbf{tepat 2 link}, kedua link
tersebut secara berpasangan harus membentuk lintasan lengkap (tanpa
putus). Kombinasi 2 link yang mungkin membentuk lintasan valid dari A ke
B adalah: 1. Link \textbf{a} bersama 1 link apa saja (Jika diasumsikan
link \emph{a} adalah bypass paralel besar) = 7 cara. 2. Link atas:
\textbf{b} seri dengan \textbf{g} = 1 cara. 3. Link tengah: \textbf{c}
seri dengan \textbf{h} = 1 cara. Total lintasan sukses dengan modal 2
link = 9 pasangan. Sementara total seluruh kombinasi memilih 2 link dari
8 link adalah \(\binom{8}{2} = 28\) kombinasi. Maka peluang bersyarat A
dapat berkomunikasi dengan B adalah:
\[ P(\text{Konek} \mid \text{Tepat 2 link hidup}) = \frac{\text{Jumlah Kombinasi Konek}}{\text{Total Kombinasi}} = \frac{9}{28} \approx \mathbf{0,3214} \]

\begin{center}\rule{0.5\linewidth}{0.5pt}\end{center}

\subsubsection{\texorpdfstring{\textbf{Soal 6: Waktu Kegagalan Komponen
(Order
Statistics)}}{Soal 6: Waktu Kegagalan Komponen (Order Statistics)}}\label{soal-6-waktu-kegagalan-komponen-order-statistics}

\textbf{Soal:} PDF Waktu \(t\) konstan 0,4 (bulan 0-2), linier ke 0
(bulan 2-3). Ada 5 komponen dioperasikan bersamaan. Tentukan peluang
\(T_5\) (waktu semua komponen rusak) melebihi 1,5 bulan!

\textbf{Jawaban:} Nilai batas waktu
\(T_5 = \max(t_1, t_2, t_3, t_4, t_5)\) akan melebihi 1,5 bulan apabila
\textbf{setidaknya satu} dari 5 komponen tersebut bertahan melewati 1,5
bulan. Hal ini setara dengan pelengkap / komplemen dari kejadian ``semua
komponen rusak sebelum atau tepat 1,5 bulan''.
\[ P(T_5 > 1,5) = 1 - P(\text{Semua } t_i \le 1,5) \] Karena tiap
komponen beroperasi independen secara statistik, probabilitas kumulatif
1 komponen (\(F(1,5)\)) dikalikan sebanyak 5 kali:
\[ P(T_5 > 1,5) = 1 - [P(t_i \le 1,5)]^5 \] Hitung \(P(t_i \le 1,5)\)
dari integral PDF yang merupakan konstanta 0,4 pada interval
\(0 \le t_0 \le 2\):
\[ P(t_i \le 1,5) = \int_{0}^{1,5} 0,4 \, dt = 0,4 \times 1,5 = 0,6 \]
Maka probabilitas \(T_5\) melebihi 1,5 bulan adalah:
\[ P(T_5 > 1,5) = 1 - (0,6)^5 = 1 - 0,07776 = \mathbf{0,92224} \]

\begin{center}\rule{0.5\linewidth}{0.5pt}\end{center}

\subsubsection{\texorpdfstring{\textbf{Soal 7: Transmisi Bit dan
Aproksimasi
Poisson}}{Soal 7: Transmisi Bit dan Aproksimasi Poisson}}\label{soal-7-transmisi-bit-dan-aproksimasi-poisson}

\textbf{Soal:} Dikirim 10 juta simbol. Peluang message \emph{error-free}
(0 error) adalah 0,5 kali dari peluang terjadinya 2 buah error. a)
Formula peluang k error (pendekatan Poisson) b) Tentukan \(p\) (peluang
error 1 simbol) dan probabilitas message \emph{error-free}!

\textbf{Jawaban:} Karena jumlah percobaan (\(n = 10.000.000\)) sangat
besar dan peluang error tiap simbol (\(p\)) sangat kecil, sebaran
binomial didekati dengan \textbf{Distribusi Poisson} dimana parameter
laju kejadian \(\lambda = n \cdot p\).

\textbf{a. Formula Peluang \(k\) Error:}
\[ P(X = k) = \frac{e^{-\lambda} \lambda^k}{k!} \]

\textbf{b. Menentukan Nilai \(p\) dan Probabilitas Bebas Kesalahan:}
Diketahui dari soal bahwa: \[ P(X = 0) = 0,5 \times P(X = 2) \]
Substitusikan ke dalam formula Poisson:
\[ \frac{e^{-\lambda} \lambda^0}{0!} = 0,5 \left( \frac{e^{-\lambda} \lambda^2}{2!} \right) \]
Karena \(e^{-\lambda}\) saling membagi (tidak nol), persamaannya dapat
disederhanakan menjadi:
\[ 1 = 0,5 \left( \frac{\lambda^2}{2} \right) \implies 1 = \frac{\lambda^2}{4} \implies \lambda^2 = 4 \implies \mathbf{\lambda = 2} \]
Karena \(\lambda = n \cdot p\), maka:
\[ 2 = 10.000.000 \times p \implies p = \frac{2}{10^7} = \mathbf{2 \times 10^{-7}} \]
Sehingga, probabilitas message benar-benar bebas dari \emph{error}
(\(X=0\)) adalah:
\[ P(X = 0) = \frac{e^{-2} 2^0}{0!} = e^{-2} \approx \mathbf{0,1353} \]




\end{document}
